% Generated from tex/chapters/ch04/09_ソフトウェア構築ツール.tex for the SWEBOK structure tree
\subsubsection{視覚的プログラミングとローコード/ゼロコードプラットフォーム}

視覚的プログラミングは、図形や部品を画面上で操作してプログラムを作る方法である。GUI ビルダーは、画面部品をドラッグアンドドロップで配置し、イベント処理のコード生成を支援する。

ローコード/ゼロコードプラットフォームは、視覚的な操作を中心にアプリケーションを作る環境である。ローコードは少量の手書きコードを必要とする。ゼロコードは、原則として手書きコードをほとんど必要としない。これらは、モデル駆動設計、視覚的プログラミング、コード生成の考え方に基づく。

このような環境は、開発速度を高める可能性がある。一方で、複雑な要件、性能、セキュリティ、拡張性、ベンダ依存、テスト容易性には注意が必要である。
